\SourceLeaf{157}
Continuing in the same way, $\Omega$ takes the form
\[
\Omega = (\nu C)^2+(\nu' C')^2+(\nu'' C'')^2+\cdots+(\nu^{(n-1)} C^{(n-1)})^2+(\beta^{(n)} B^{(n)})^2,
\]
where all the $C$ are independent of $x$.



\Needspace{6\baselineskip}
\begin{center}
\textsc{Gauss} to \textsc{Olbers}. Göttingen, 14 April 1819.
\end{center}

\begin{center}
. . . . . . . . . . . . . . . .
\end{center}

Likewise, I do not quite understand your objection to the correct formula for the weight of the result of combining 2 determinations with weights $p$, $q$, namely $\dfrac{pq}{p+q}$, unless you did not compare my letter with the passage in \textsc{Littrow}. I do indeed assume that the partial results are \textls{independent} of one another. But the issue is \emph{what kind of} combination is meant. In the \emph{present case}, the fact that $\dfrac{pq}{p+q}$ is always $<p$ and $<q$ is quite reasonable. The opposite would be absurd. Take another very similar case. \textsc{Schumacher} intends to measure an arc of longitude. Let the endpoints be $A$ and $C$. He cannot compare them, that is, his rockets, directly, and needs an intermediate point $B$. If he first compares $A$ with $B$ 100 times, and then $C$ with $B$ also 100 times, the partial differences in longitude $B-A$, $C-B$ have a high and equal reliability, whose weight is given by $100=p$, $100=q$. But the difference between the locations $C$, $A$ derived by combining the two will clearly be less reliable than if it had been derived from 100 direct comparisons. Only probability theory can tell how much lower the weight of the result of this combination is. It shows that this \textls{inferred} result is only as accurate as the mean of 50 direct comparisons. Only if $q=\infty$, and hence $C-B$ is absolutely exact, will the \textls{inferred} result for $C-A$ be as accurate as the $B-A$ used, as indeed follows from the nature of the case. \textsc{Littrow}'s incorrect formula would give for $q=\infty$ the absurd result $4p$. --- I suppose that \textsc{Littrow} was led astray by the fact that the
